\documentclass[11pt,a4paper]{article}

\usepackage[utf8]{inputenc}
\usepackage[T1]{fontenc}
\usepackage{amsmath}
\usepackage{amssymb}
\usepackage{graphicx}
\usepackage{hyperref}
\usepackage{booktabs}
\usepackage{xcolor}
\usepackage[authoryear,round]{natbib}
\usepackage[margin=1in]{geometry}

\title{MedScreen: Evidence-Based Factuality Filtering of Medical Training Data}

\author{Ciaran O'Brien}

\date{\today}

\begin{document}

\maketitle

\begin{abstract}

Medical language models inherit the errors of the corpus they train on. A published paper can be fluent, well cited, and formally retracted, or state a claim that later evidence overturned, and a filter that judges papers by how they read will keep it. MedScreen is an evidence-based factuality filter for medical training data. It extracts the specific claims a PubMed paper makes, retrieves trusted studies that bear on each claim, judges whether that evidence supports or refutes it, and aggregates the result into a per-paper verdict and a recommended training action: keep, down-weight, drop, or flag for review. Because the filter is only as good as its search, we test that one dependency directly: on thirty-two consensus reversals whose disproving studies are known in advance, the search finds the disproving study for 94\% of cases (30 of 32) without any model. With a lightweight judge (Gemini 2.5 Flash Lite), 91\% of those studies are then read as refuting, and no true control paper is ever wrongly dropped (0 of 32), because a drop requires two independent strong studies to agree, so a single misjudged control is down-weighted rather than deleted. The system judges factual correctness alone and never lets writing quality, journal prestige, or author reputation stand in for truth. We report where the search fails, why judging abstracts rather than full text limits precision, and what a production version would need.

\end{abstract}

\section{Introduction}

Language models trained on medical text reproduce what that text asserts. If the training corpus contains claims that are false, superseded, or drawn from retracted work, the model learns them and repeats them with confidence. Data quality therefore sets a ceiling on model trustworthiness, and for medical models the cost of a confidently wrong output is high.

The usual ways of screening a corpus judge a paper by its surface. Rule-based filters key on formatting and vocabulary; learned quality raters score fluency and structure; citation and journal-prestige signals reward reputation. Confident medical misinformation imitates all of these well. A retracted paper reads like a sound one, a reversed claim was mainstream and heavily cited for years before the evidence turned, and neither is distinguishable from good work by surface features alone.

MedScreen takes a different axis. It asks one question of each paper: are its claims factually correct given trusted scientific evidence. It answers by retrieving that evidence and checking the claims against it, rather than inferring truth from how the paper reads. This is deliberately a factuality filter and nothing else. It does not judge sample size, statistical rigour, reproducibility, writing quality, journal prestige, or author h-index, and it never lets those stand in for truth. Paper quality is assumed to be learned from the corpus itself; the filter's only job is to remove provably refuted claims and flag ungrounded ones.

The central risk in this design is retrieval. If the study that contradicts a wrong claim is never retrieved, no later step can use it, and the paper passes. So alongside the filter we build a validation harness that measures exactly this dependency on claims the field already knows were reversed, where the disproving study is recorded in advance. The filter is the product; the harness is the scaffold that shows the product can work.

This paper describes the filter, the retrieval and scoring it rests on, and the validation results from a proof-of-concept implementation. We report both what works (the search finds the disproving study 94\% of the time, using no model, and never wrongly drops a control paper) and where it does not (two searches that miss a conceptual reversal, and a limit on the judge from reading abstracts rather than full text).

\section{Related Work}

The closest methodological cousins decompose text into atomic claims and verify each against retrieved evidence. FActScore \citep{factscore2023} breaks a generation into atomic facts and scores each as supported or unsupported against a knowledge source. SAFE \citep{safe2024} extends this with a search-augmented evaluator that issues queries and rates each atomic fact against retrieved results. FEVER \citep{fever2018} established the retrieve-then-verify task and dataset for open-domain claims.

All three apply the pipeline to model \emph{outputs}: they check what a system generated. MedScreen applies the same retrieve-then-verify decomposition to training \emph{inputs}, the published papers that make up a corpus, and adds two things suited to that setting. First, an evidence-tier weighting, so a randomised trial moves a verdict more than a case report. Second, a retrieval-recall validation harness anchored on known consensus reversals, which measures the one dependency the filter's accuracy rests on against ground truth.

Reputation-based screening is the implicit alternative and the one MedScreen argues against. Citation count and journal prestige judge who is speaking, not whether they are right. The belief that hormone replacement therapy prevents coronary heart disease was heavily cited the entire time it was wrong, until the Women's Health Initiative trial \citep{whi2002} found the opposite.

\section{Method}

For each paper the filter runs five stages. An LLM is used in only two of them, claim extraction and stance judgment, and never runs retrieval or decides the keep or drop verdict. The whole pipeline also runs offline on deterministic stub backends so the control flow can be checked without a key or a network call.

\subsection{Ingest}

Input is PubMed/MEDLINE XML rather than plain text, because the XML carries the signals the filter relies on: the \texttt{CommentsCorrectionsList} retraction and comment links, the publication types that set evidence tiers, and the MeSH terms. The ingester parses each file into a structured record and skips, with an error, any file it cannot read or that is not a PubMed article set.

\subsection{Claim extraction}

An LLM lifts the paper's specific claims out of its text, keeping the conditions attached: population, comparator, dose, direction, and setting. Condition-stripped claims manufacture false verdicts, so retaining them is treated as an invariant rather than a nicety.

\subsection{Evidence retrieval}

For each claim the filter builds several keyword and boolean queries and sends them to the PubMed E-utilities and Europe PMC search endpoints, which return capped, relevance-ranked lists of study IDs. It unions those with any dispute links from the paper's own XML, then fetches each candidate's abstract and publication type. Evidence is not stored locally and this is not a vector search over all of PubMed.

Retrieval is tuned so a known refutation actually comes back. Several searches run per claim and are unioned: a broad one on intervention and outcome, one limited to strong study types, one targeting contradiction terms such as \emph{risk}, \emph{harm}, and \emph{no benefit}, a retraction-targeted one so link expansion can reach a retraction notice, and a condition-focused one on intervention and population, since a descriptive outcome over-narrows while the disease name pins a landmark trial. Two sources are queried independently. Embeddings are used only to re-rank an already-fetched pool during validation; they cannot recover a study the queries never returned. Whether the queries succeed is the central failure mode, and exactly what the validation measures.

\subsection{Stance judgment}

An LLM labels each candidate study as supporting, refuting, or neutral toward the claim, and reports a confidence with each label. The judge reads only each study's title and abstract, never its full text. This is a deliberate cost trade-off: abstracts are cheap to fetch and short to send, but a refutation or a condition caveat that lives only in a paper's Results or Methods is invisible to the judge. It is a known limit on how well the judge can do, accepted for the proof of concept.

\subsection{Scoring and verdict}

Scoring is a mechanical aggregation of the stance labels, weighted by evidence tier, with fixed thresholds. The arithmetic and the keep or drop decision are deterministic and never call the model, but the labels and confidences they combine are model judgments, so the stance judge is load-bearing for the direction of a verdict even though it never makes the verdict itself.

A retraction fast path runs first. If the paper's own XML shows it is formally retracted, via a \texttt{RetractionIn} link or the \texttt{Retracted Publication} publication type, it is dropped immediately with \texttt{verdict\_basis = retraction}, skipping extraction, retrieval, and stance. This is the cheapest and strongest signal. A present retraction signal is reliable; its absence is not proof a paper was not retracted, given indexing lag, so every other paper is scored on retrieved evidence.

Each retrieved study is weighed by its publication-type evidence tier (guideline 1.0, retraction 0.95, systematic review 0.9, meta-analysis 0.85, RCT 0.8, observational 0.5, case report 0.2, anything else 0.4). Its pull on a claim is that tier times the stance judge's confidence. Refutations the judge marks off-scope (a different population, dose, comparator, or setting than the claim asserts) are discarded first, since such a study refutes a different claim and is not evidence against this one. A claim's score starts at 0.5, rises with aggregated supporting evidence and falls with aggregated refuting evidence, with refutation weighed roughly twice as heavily. The aggregate is not the single strongest study: the strongest sets the base and each further agreeing study adds a diminishing share, so a consistent body of evidence counts for more than one study while many weak studies cannot overpower one strong study, because each study's weight already reflects its tier.

The per-claim verdict is one of \texttt{refuted}, \texttt{contested}, \texttt{supported}, \texttt{neutral} (only neutral evidence found), or \texttt{ungrounded} (no evidence found). A claim is \texttt{refuted}, which drops its paper, only when the refutation is strong and corroborated: at least two separate studies refute it, the combined refuting strength is at least 0.6, and the strongest of those studies is tier 0.8 or higher with a stance confidence of at least 0.7. Requiring two studies means a single refutation of one claim in a multi-claim paper down-weights rather than drops it, which was the main cause of wrongful drops. Anything weaker, or any claim with evidence on both sides, is \texttt{contested}.

The paper takes the verdict and action of its single most damning claim: its score is the lowest claim score, its verdict the worst. The mapping to an action is fixed, shown in Table~\ref{tab:actions}. A paper correct when written and later superseded but never actually contradicted is kept on purpose, because a once-true paper is not false and science is incremental. Only a genuine subsequent refutation moves a paper off \texttt{keep}.

\begin{table}[h]
\centering
\begin{tabular}{lll}
\toprule
\textbf{Verdict} & \textbf{Action} & \textbf{Meaning} \\
\midrule
\texttt{refuted} & drop & High-tier evidence contradicts the claim. \\
\texttt{contested} & downweight & Evidence points both ways, or the refutation is weak. \\
\texttt{supported} & keep & Evidence backs the claim, no credible refutation. \\
\texttt{neutral} & keep & Only neutral evidence found; a missing refutation is not disproof. \\
\texttt{ungrounded} & review & No evidence found at all; flagged for a human. \\
\bottomrule
\end{tabular}
\caption{Verdict-to-action mapping. A paper inherits the action of its most damning claim.}
\label{tab:actions}
\end{table}

\section{Validation}

The filter is only as good as its search, so the validation harness measures retrieval directly, on claims the field already knows were wrong and whose disproving study is recorded in advance. It runs the filter's search and checks how often the known study comes back, scored in stages so a failure traces to the right one.

\subsection{Gold set}

The gold slice holds thirty-two reversed claims and thirty-two controls. The reversed claims split into two kinds. A \emph{reversal} is good-faith science superseded by a newer study, found by keyword and high-tier search: twenty-eight such cases, including HRT and coronary heart disease overturned by WHI \citep{whi2002}, antiarrhythmics after myocardial infarction overturned by CAST \citep{cast1989}, and vertebroplasty for osteoporotic fractures overturned by a sham-controlled trial \citep{buchbinder2009}. A \emph{fabrication} is retracted misconduct whose disproving evidence is the retraction notice, reached by a retraction-targeted query: four such cases. The thirty-two controls are still-supported claims that must be kept. Every disproving PMID was verified against PubMed esummary; the gold set is the most accuracy-critical artifact in the project.

\subsection{Metrics}

We report four measures. The first, whether the search brought the known disproving study back at all, does not use the model and is the primary number. The second, ranking, asks whether that study also landed near the top of the retrieved list, since only the twenty most relevant per claim are shown to the judge. The third, whether the judge then reads a disproving study as refuting, isolates the judge from the search. The fourth, how often a still-true control picks up any refuting label, guards against a filter that flags everything and so is worthless. Every miss is tagged with its cause, from ``never retrieved'' to ``retrieved but not recognised'', so a failure points at the stage that caused it.

\subsection{Results}

Results are summarised in Table~\ref{tab:results}. The search found the disproving study for 94\% of the reversed claims, measured without the model. When the judge was then shown that study, it read it as refuting 91\% of the time. This 91\% always shows the disproving study to the judge; the figure that reflects what a live filter actually forwards (its twenty most relevant studies per claim) is 78\%, so the real limit is which studies rank high enough to be judged, not the judge itself. On the still-true controls, 47\% picked up at least one refuting label, because the controls are ordinary claims whose retrieved evidence carries legitimate caveats that an abstract-only judge over-reads as refutation. The judge itself marks most of these as off-scope (a different population, dose, comparator, or setting than the claim), and the scorer discards an off-scope refutation as evidence about a different claim, so the rate that reflects what the scorer actually counts falls to 25\% (8 of 32). Either way not one control was dropped (0 of 32): a drop requires two independent strong studies to agree, so each over-flagged control was down-weighted, a reversible action, rather than deleted.

\begin{table}[h]
\centering
\begin{tabular}{lr}
\toprule
\textbf{Metric} & \textbf{Result} \\
\midrule
Disproving study found by the search & 94\% (30/32) \\
Judge reads it as refuting, once shown & 91\% (29/32) \\
True controls flagged, scope-aware & 25\% (8/32) \\
True controls picking up any refuting label (raw) & 47\% (15/32) \\
True controls wrongly dropped & 0/32 \\
Claims correctly extracted & 83\% \\
Claim conditions kept intact & 93--100\% \\
\bottomrule
\end{tabular}
\caption{Validation results on the full gold slice: 28 reversals, 4 fabrications, 32 controls. The first row uses no model; the stance and control rows use Gemini 2.5 Flash Lite.}
\label{tab:results}
\end{table}

Claim extraction, comparing Gemini 2.5 Flash Lite against a stronger reference model, found 83\% of the expected claims and kept their conditions 93 to 100\% of the time. Its precision is lower, near 43\%, because it extracts more, finer-grained claims than the reference. The primary result, whether the search finds the disproving study, uses no model and so can be checked at no cost; only the judge and control measures need a paid model.

\section{Analysis}

\subsection{The two retrieval misses}

Two reversals are not recovered, and both are accepted limitations rather than bugs. The peptic-ulcer etiology reversal shares only the disease with its refutation: the claim is about stress and acid, the answer is \emph{Helicobacter pylori} \citep{marshall1984}, so no keyword or MeSH query reaches it without already naming the bacterium. This is a known limit of keyword search for conceptual reversals. The vertebroplasty reversal's landmark trial is buried among many similar high-tier trials that the condition alone does not disambiguate. Bridging either would need a semantic query-expansion step, which is deferred on purpose rather than risking control precision by over-broadening the search. The filter is precision first, so some misses are acceptable.

\subsection{Behaviour on ordinary papers}

The harness measures known reversals, which are the easiest cases. To see behaviour on ordinary papers, the filter was run on thirty recent, non-retracted papers across common clinical areas. It kept fifteen, down-weighted fourteen, and dropped one, so half were not kept. That rate is high, driven mainly by review and guideline papers: they make many claims and so are more exposed to a single flagged claim than a primary research paper, under the most-damning-claim rollup. The single drop was a corroborated refutation of one claim, which is exactly what the drop action is reserved for.

\subsection{The abstract-only ceiling}

The stance judge reads titles and abstracts, not full text. A refutation or a caveat stated only in a paper's Results or Methods cannot be seen, which limits how well the judge does. The over-flagging on the controls is this judge limitation, not a search one: the search surfaced the right evidence, and the judgment on the abstract alone was ambiguous.

\section{Limitations and Future Work}

Several limits are deliberate for a proof of concept. Evidence is weighted by a flat publication-type tier, which rates a study by its type label alone; GRADE \citep{grade2008} would keep the same role but adjust for risk of bias, inconsistency, indirectness, and imprecision, so evidence quality rather than study type drives the verdict. Retrieval uses brute-force cosine over a fetched pool; an approximate-nearest-neighbour index would scale it. Query construction is keyword and boolean; LLM-driven query and topic expansion would match on meaning and reach the conceptual reversals that keyword search cannot. Stance reads abstracts only; reading full text would give the judge the context to rule out false refutations. All of these are on the roadmap and none is built.

The evaluation has its own gaps. The results come from sixty-four claims (twenty-eight reversals, four fabrications, and thirty-two controls), a mix of famous and less-famous cases so the search is not judged only on the easiest ones. The search finds the disproving study 94\% of the time and never drops a control, but 47\% of controls pick up a spurious refuting label, which shows ordinary controls exercise the judge harder than famous reversals do. There is still no end-to-end accuracy figure on a human-labelled set of ordinary papers, which is the most informative evaluation to add. The extraction scoring run against its reference set is written but deferred, because running the extractor needs a real LLM. The reference claims themselves were written by a strong model rather than a person, so the extraction check measures agreement between two models, not agreement with human ground truth.

\section{Conclusion}

A medical training corpus can be screened for factuality by retrieving evidence and checking each paper's claims against it, rather than judging the paper by how it reads. MedScreen does this with an LLM confined to extraction and stance, a deterministic tier-weighted scorer that makes the keep or drop decision, and a retraction fast path for the strongest signal. Because the whole approach rests on the search, we measure the search directly: on a set of known reversals, it recovers 94\% of disproving studies with no model, and the strict corroboration thresholds mean no true control is wrongly dropped even when the judge is imperfect. The result is a transparent, auditable filter whose verdicts show the evidence they rest on. Before production use, stance judgment and query precision are the two things to improve, and the paths to both are clear.

\section*{Code and Data Availability}

The MedScreen proof-of-concept implementation, including the validation harness and the gold dataset of consensus reversals, is publicly available at \url{https://github.com/obrienciaran/MedScreen_filter_POC} \citep{medscreen}.

\bibliographystyle{plainnat}
\bibliography{medscreen}

\end{document}
